\begin{tikzpicture}[
  font=\scriptsize,
  execute at begin node={\hyphenpenalty=10000\exhyphenpenalty=10000},
  label/.style={text width=2.25cm, minimum height=0.95cm, align=left},
  pengguna/.style={draw, rounded corners, fill=praditagreen!15, text width=2.6cm, minimum height=0.95cm, align=center},
  pengalaman/.style={draw, rounded corners, fill=orange!12, text width=2.6cm, minimum height=0.95cm, align=center},
  proses/.style={draw, rounded corners, fill=blue!8, text width=2.6cm, minimum height=0.95cm, align=center},
  sistem/.style={draw, rounded corners, fill=blue!8, text width=2.6cm, minimum height=0.95cm, align=center},
  sumber/.style={draw, rounded corners, fill=gray!35, text width=2.6cm, minimum height=0.95cm, align=center},
  panah/.style={-{Stealth[length=2mm]}, thick},
  lewat/.style={panah, dashed}
]
  % Kolom kiri menyebut alasan tiap lapisan berubah, karena di situlah nilainya.
  \node[label] at (0,0) {\textbf{Pengguna}};
  \node[label] at (0,-1.45) {\textbf{Pengalaman}\\{\tiny\mdseries berubah mengikuti kanal}};
  \node[label] at (0,-2.9) {\textbf{Proses}\\{\tiny\mdseries berubah mengikuti aturan bisnis}};
  \node[label] at (0,-4.35) {\textbf{Sistem}\\{\tiny\mdseries berubah mengikuti vendor}};
  \node[label] at (0,-5.8) {\textbf{Sumber}};

  \node[pengguna] (u1) at (3.0,0) {Medikara+ dan 40 klinik mitra};
  \node[pengguna] (u2) at (6.0,0) {3 laboratorium dan 40 klinik mitra};
  \node[pengguna] (u3) at (9.0,0) {9 asuransi};

  \node[pengalaman] (e1) at (3.0,-1.45) {Jadwal, lokasi, dan status pada satu tampilan};

  \node[proses] (p1) at (3.0,-2.9) {Pemesanan: cari, pilih, kunci slot};
  \node[proses] (p2) at (6.0,-2.9) {Pengiriman hasil: tulis, beri tahu, ambil};
  \node[proses] (p3) at (9.0,-2.9) {Klaim: susun, periksa, kirim};

  \node[sistem] (s1) at (3.0,-4.35) {Adaptor jadwal untuk 5 SIMRS};
  \node[sistem] (s2) at (6.0,-4.35) {Adaptor hasil untuk SIMRS};
  \node[sistem] (s3) at (9.0,-4.35) {Adaptor tagihan untuk SIMRS dan ERP};

  \node[sumber] (b1) at (3.0,-5.8) {5 SIMRS\\{\tiny\mdseries struktur berbeda}};
  \node[sumber] (b2) at (6.0,-5.8) {Laboratorium dan SIMRS};
  \node[sumber] (b3) at (9.0,-5.8) {5 SIMRS dan ERP keuangan};

  \draw[panah] (u1) -- (e1);
  \draw[panah] (e1) -- (p1);
  \draw[lewat] (u2) -- (p2);
  \draw[lewat] (u3) -- (p3);
  \foreach \kolom in {1,2,3} {
    \draw[panah] (p\kolom) -- (s\kolom);
    \draw[panah] (s\kolom) -- (b\kolom);
  }

  \node[font=\tiny, align=center, text width=8.6cm] at (6.0,-6.6) {Panah putus-putus: lapisan pengalaman dilewati karena tidak menambah nilai bagi pemakai ini};
\end{tikzpicture}
